\section{Conclusions}
\label{sec:conclusions}

This study developed an offline, grid-based framework for diagnosing the
feasibility and geometry of trans-Arctic shipping routes under static
bathymetric and historical sea-ice constraints.  Using GEBCO~2024
bathymetry regridded to a $0.5^{\circ}$ pan-Arctic corridor
($40$--$85^{\circ}$N, $20^{\circ}$W–$180^{\circ}$E), we constructed a
sea-only routing mask and applied A* search to a canonical west–east
origin–destination pair.  Even without environmental forcing, enforcing
strict sea-only feasibility increased route length from $4{,}150$~NM
(great-circle) to about $4{,}560$~NM (sea-only), i.e.\ a $\sim 10\%$
penalty arising purely from the shape of coastlines and islands on a
realistic bathymetric grid.

Static depth constraints at $0.5^{\circ}$ resolution were found to be
relatively benign at basin scale.  For representative under-keel limits
of $h_{\min}=20$–$50$~m, more than $85\%$ of the sea mask remained
depth-safe, and a sea-only route persisted for $h_{\min}=20$~m.
Depth therefore acts as a structural constraint---eliminating shallow
shelves and passages---but does not, by itself, prevent large-scale
trans-Arctic connectivity in this configuration.  Bathymetry still
matters, however, for shaping where such routes can run and for
interpreting how close candidate paths approach shallow and coastal
zones.

Historical sea-ice concentration emerged as the dominant
time-varying control.  Using summer 2018 CMEMS sea-ice concentration
regridded to the same $0.5^{\circ}$ routing grid and an ice-safety
threshold of $\mathrm{SIC}<15\%$, the fraction of ice-safe sea increased
from roughly $60\%$ in mid-June to over $80\%$ in mid-September.
Continuous ice-safe trans-Arctic routes appeared only from mid-August
onwards and were consistently longer than the sea-only benchmark:
$\sim 5{,}200$~NM in mid-August, decreasing to $\sim 5{,}000$~NM by
late September as the marginal ice zone retreated.  The excess-distance
penalty relative to great-circle therefore drops from about $25\%$ to
$20\%$ over this late-summer window, highlighting how the timing of
Arctic transits trades off against both feasibility and path length.

At the large scales considered here, route geometry was only weakly
sensitive to the precise sea-ice threshold once a route existed.
Varying the concentration threshold between $5$ and $50\%$ changed the
fraction of ice-safe sea by a few percentage points and altered
route length by only 20–30 NM.  For coarse-grid,
basin-scale diagnostics this suggests that the choice of SIC threshold
is less critical than the choice of season and year, although the
threshold will become more important once narrower passages and coastal
traffic lanes are resolved.

Combining bathymetry and ice into a joint feasibility mask exposed a
more restrictive picture.  For a relatively mild depth constraint
($h_{\min}=20$~m) and $\mathrm{SIC}<15\%$ on 15~September~2018, only
about $75\%$ of sea cells were simultaneously depth-safe and ice-safe,
split into 35 disconnected components.  In this joint mask the canonical
origin and destination fall into different components, and no fully
depth- and ice-safe trans-Arctic route exists at $0.5^{\circ}$
resolution.  This underlines that optimistic statements about a
``seasonally ice-free'' Arctic can overstate actual navigable
connectivity once realistic draft constraints are imposed.

The present framework has several limitations that also point directly
to future work.  First, the routing grid is coarse
($0.5^{\circ}\approx 55$~km), so narrow straits, detailed coastline
geometry and port approaches are not resolved; the results are intended
as basin-scale diagnostics, not operational tracks.  Second, only a
single canonical origin–destination pair and a single summer season
(2018) were analysed; extending the analysis to a portfolio of
Europe–Asia and intra-Arctic ODs and to multi-year sea-ice statistics
would better capture inter-annual variability.  Third, cost was defined
purely as geographic distance, without explicit representation of
metocean forcing, resistance, fuel consumption, emissions, or risk.
Finally, several safety-relevant factors---sea-ice thickness, ridging,
traffic separation schemes, and region-specific regulatory corridors---are
omitted and should be incorporated before any operational use.

Despite these limitations, the study provides a physically grounded
baseline for large-scale, safety-constrained Arctic routing.  It
quantifies how static bathymetry and historical sea ice jointly shape
the feasibility, timing and excess distance of trans-Arctic connections,
and it delineates when such routes are even theoretically open.
Follow-on work will build directly on this foundation by (i) increasing
spatial resolution and adding multiple ODs, (ii) integrating
time-evolving sea-ice forecasts and metocean forcing into dynamic
cost fields, and (iii) extending the single-distance metric towards
multi-objective formulations that link routing more directly to
fuel consumption, emissions and schedule reliability.